\begin{lemma}
\label{lemma-filter-quasi-compact}
Let $S$ be a scheme. Let $X$ be a quasi-compact algebraic space
over $S$. There exist open subspaces
$$
\ldots \subset U_4 \subset U_3 \subset U_2 \subset U_1 = X
$$
with the following properties:
\begin{enumerate}
\item setting $T_p = U_p \setminus U_{p + 1}$ (with reduced induced subspace
structure) there exists a separated scheme $V_p$ and a surjective \'etale
morphism $f_p : V_p \to U_p$ such that $f_p^{-1}(T_p) \to T_p$ is an
isomorphism,
\item if $x \in |X|$ can be represented by a quasi-compact morphism
$\Spec(k) \to X$ from a field, then $x \in T_p$ for some $p$.
\end{enumerate}
\end{lemma}

\begin{proof}
By Properties of Spaces, Lemma
\ref{spaces-properties-lemma-quasi-compact-affine-cover}
we can choose an affine scheme $U$ and a surjective \'etale morphism
$U \to X$. For $p \geq 0$ set
$$
W_p = U \times_X \ldots \times_X U \setminus \text{all diagonals}
$$
where the fibre product has $p$ factors. Since $U$ is separated,
the morphism $U \to X$ is separated and all fibre products
$U \times_X \ldots \times_X U$ are separated schemes. Since $U \to X$ is
separated the diagonal $U \to U \times_X U$ is a closed immersion. Since
$U \to X$ is \'etale the diagonal $U \to U \times_X U$ is an open
immersion, see Morphisms of Spaces, Lemmas
\ref{spaces-morphisms-lemma-etale-unramified} and
\ref{spaces-morphisms-lemma-diagonal-unramified-morphism}.
Similarly, all the diagonal morphisms are open and closed immersions and
$W_p$ is an open and closed subscheme of $U \times_X \ldots \times_X U$.
Moreover, the morphism
$$
U \times_X \ldots \times_X U \longrightarrow
U \times_{\Spec(\mathbf{Z})} \ldots \times_{\Spec(\mathbf{Z})} U
$$
is locally quasi-finite and separated (Morphisms of Spaces,
Lemma \ref{spaces-morphisms-lemma-fibre-product-after-map})
and its target is an affine scheme. Hence every finite set of points of
$U \times_X \ldots \times_X U$ is contained in an affine open, see
More on Morphisms, Lemma
\ref{more-morphisms-lemma-separated-locally-quasi-finite-over-affine}.
Therefore, the same is true for $W_p$.
There is a free action of the symmetric group $S_p$ on $W_p$ over $X$
(because we threw out the fix point locus from
$U \times_X \ldots \times_X U$). By the above and
Properties of Spaces, Proposition
\ref{spaces-properties-proposition-finite-flat-equivalence-global}
the quotient $V_p = W_p/S_p$ is a scheme. Since the action of
$S_p$ on $W_p$ was over $X$, there is a morphism $V_p \to X$.
Since $W_p \to X$ is \'etale and since $W_p \to V_p$ is surjective
\'etale, it follows that also $V_p \to X$ is \'etale, see
Properties of Spaces, Lemma \ref{spaces-properties-lemma-etale-local}.
Observe that $V_p$ is a separated scheme by
Properties of Spaces, Lemma
\ref{spaces-properties-lemma-quotient-separated}.

\medskip\noindent
We let $U_p \subset X$ be the open subspace which is the
image of $V_p \to X$. By construction a morphism $\Spec(k) \to X$ with
$k$ algebraically closed, factors through $U_p$ if and only if
$U \times_X \Spec(k)$ has $\geq p$ points; as usual observe that
$U \times_X \Spec(k)$ is scheme theoretically a disjoint union of
(possibly infinitely many) copies of $\Spec(k)$, see
Remark \ref{remark-recall}. It follows that
the $U_p$ give a filtration of $X$ as stated in the lemma.
Moreover, our morphism $\Spec(k) \to X$ factors through $T_p$
if and only if $U \times_X \Spec(k)$ has exactly $p$ points.
In this case we see that $V_p \times_X \Spec(k)$ has exactly one point.
Set $Z_p = f_p^{-1}(T_p) \subset V_p$. This is a closed subscheme of $V_p$.
Then $Z_p \to T_p$ is an \'etale morphism between
algebraic spaces which induces a bijection on $k$-valued
points for any algebraically closed field $k$. To be sure this
implies that $Z_p \to T_p$ is universally injective, whence an
open immersion by
Morphisms of Spaces, Lemma
\ref{spaces-morphisms-lemma-etale-universally-injective-open}
hence an isomorphism and (1) has been proved.

\medskip\noindent
Let $x : \Spec(k) \to X$ be a quasi-compact morphism where $k$ is a field.
Then the composition $\Spec(\overline{k}) \to \Spec(k) \to X$ is quasi-compact
as well (Morphisms of Spaces, Lemma
\ref{spaces-morphisms-lemma-composition-quasi-compact}).
In this case the scheme $U \times_X \Spec(\overline{k})$ is
quasi-compact. In view of the fact (seen above) that it is a disjoint union
of copies of $\Spec(\overline{k})$ we find that it has finitely many points.
If the number of points is $p$, then we see that indeed $x \in T_p$ and
the proof is finished.
\end{proof}
